\documentclass{article}
\usepackage{amsmath,amssymb,xurl}
\usepackage[hidelinks]{hyperref}
% Shared commands for notes in docs/paper-gaps/.

\DeclareUnicodeCharacter{2080}{\ifmmode{}_{0}\else\textsubscript{0}\fi}
\DeclareUnicodeCharacter{2081}{\ifmmode{}_{1}\else\textsubscript{1}\fi}
\DeclareUnicodeCharacter{2082}{\ifmmode{}_{2}\else\textsubscript{2}\fi}
\DeclareUnicodeCharacter{2099}{\ifmmode{}_{n}\else\textsubscript{n}\fi}

\newcommand{\Lean}{\textsc{Lean}}
\ExplSyntaxOn
\NewDocumentCommand{\leanid}{m}{
  \ifmmode
    \textsf{\detokenize{#1}}
  \else
    \group_begin:
    \urlstyle{sf}
    \tl_set:Nn \l_tmpa_tl {#1}
    \tl_replace_all:Nnn \l_tmpa_tl {\_} {_}
    \exp_args:NV \path \l_tmpa_tl
    \group_end:
  \fi
}
\ExplSyntaxOff
\providecommand{\tr}{\operatorname{tr}}
\newcommand{\tnleanIssueBase}{https://github.com/LionSR/TNLean/issues/}
\newcommand{\tnleanPrBase}{https://github.com/LionSR/TNLean/pull/}
\newcommand{\tnleanDocBase}{https://sirui-lu.com/TNLean/docs/}
\newcommand{\ghissue}[1]{\href{\tnleanIssueBase#1}{GitHub issue \##1}}
\newcommand{\ghpr}[1]{\href{\tnleanPrBase#1}{PR \##1}}
\newcommand{\doclink}[1]{\href{\tnleanDocBase#1.html}{\path{#1}}}

% Dirac notation and the identity operator, as in blueprint/src/macros/common.tex.
% \providecommand keeps note-local definitions authoritative.
\usepackage{dsfont}
\providecommand{\ket}[1]{|#1\rangle}
\providecommand{\bra}[1]{\langle #1|}
\providecommand{\braket}[2]{\langle #1 | #2 \rangle}
\providecommand{\ketbra}[2]{|#1\rangle\!\langle #2|}
\providecommand{\Id}{\mathds{1}}
\providecommand{\one}{\mathds{1}}
\providecommand{\C}{\mathbb{C}}
\providecommand{\MN}[1]{M_{#1}(\C)}
\providecommand{\MD}{M_D(\C)}

% Verdict marker: \gapnote{<kind>}{<status>}. Kind is the policy
% classification (clarification, local-correction, scope-restriction,
% unfaithful, false-source, open-gap); status is open, wip (elimination
% underway), resolved, or historical. Machine-read by the site index;
% prints a verdict line.
\newcommand{\gapnote}[2]{\par\noindent\textbf{Verdict:} #1 (#2).\par}

\title{An Explicit Endpoint-Free Native Annular Flux Observation}
\author{}
\date{}
\begin{document}
\maketitle
\gapnote{scope-restriction}{open}

Schuch, Cirac, and P\'erez-Garc\'ia,
\href{https://arxiv.org/abs/1001.3807v3}{arXiv:1001.3807v3},
Lemma 6.14, source \path{Papers/1001.3807/paper_v3.tex}, lines 2199--2214,
compares local physics away from flux endpoints with the vacuum. This
cannot hold for arbitrary endpoint-free annular observations.

The theorem
\leanid{TNLean.PEPS.IsGIsometric.torusAnnularFlux_normalized_density_ne_vacuum}
constructs an explicit obstruction in the seven-by-seven torus. The
observation consists of the twelve physical perimeter sites of the square
from zero to three in both coordinates. Its four interior sites are
unobserved. The actual native dual path consists of four east steps from
plaquette $(1,1)$ to plaquette $(5,1)$. Every physical corner of both
endpoint plaquettes is outside the observed ring, as proved by
\leanid{TNLean.PEPS.torusAnnularFluxPath_endpoint_corners_outside}.

For any finite group and any nonidentity insertion element, the actual
perimeter holonomy is the inverse insertion element. This is computed
from the native path's signed crossings and the actual ordered-edge
transport; no nontrivial holonomy is assumed. The observed perimeter
graph is derived connected. The complete physical closed-walk measurement
then distinguishes this state from vacuum for every native regular
group-isometric tensor. Both coefficient cuts are proved nonzero and their
trace-normalized observed row Grams are unequal. No density inequality,
measurement operator or nonzero witness is supplied.

This is an explicit counterexample to endpoint-only vacuum-density wording,
not a classification of all annular observations. It leaves unchanged the
two simple connectedness hypotheses and endpoint exclusions in the positive
native theorem, and the envelope hypotheses in its subregion extension.
It is also distinct from the nonabelian test-group obstruction to restricted
homotopy. The broader admissible geometry remains open under
\ghissue{8703}; \ghissue{8947} is unchanged.
\end{document}
